\section{AHSME 1961}
        \begin{enumerate}
        	\item (\textit{AHSME 1961}) When simplified $(-\frac{1}{125})^{-2/3}$ becomes:


 $\textbf{(A)}\ \frac{1}{25} \qquad \textbf{(B)}\ -\frac{1}{25} \qquad \textbf{(C)}\ 25\qquad \textbf{(D)}\ -25\qquad \textbf{(E)}\ 25\sqrt{-1}$

 \textbf{Answer: }C



\textbf{Solution 1}

To remove the negative exponent, flip the fraction of the base.  This results in $(-125)^{2/3}$.

 Then, cube root $-125$ and and square the result to get the answer of $25$, or answer choice $\boxed{\textbf{(C)}}$.

 



\textbf{Solution 2}

<i>By Ajeet Dubey | \href{https://ioqm.in}{ioqm.in}</i>

 \textbf{Step 1: Rewrite the negative exponent.}

 For any nonzero number $u$, we have $u^{-2/3}=1/u^{2/3}$. Hence$\left(-\frac1{125}\right)^{-2/3}=\frac{1}{\left(-\frac1{125}\right)^{2/3}}.$

 \textbf{Step 2: Evaluate the fractional power.}

 The real cube root of $-1/125$ is $-1/5$, so$\left(-\frac1{125}\right)^{2/3}=\left(-\frac15\right)^2=\frac1{25}.$Therefore the required value is $25$.

 Thus the answer is $\boxed{\textbf{(C)}}$.

 


	\item (\textit{AHSME 1961}) An automobile travels $a/6$ feet in $r$ seconds. If this rate is maintained for $3$ minutes, how many yards does it travel in $3$ minutes?


 $\textbf{(A)}\ \frac{a}{1080r}\qquad \textbf{(B)}\ \frac{30r}{a}\qquad \textbf{(C)}\ \frac{30a}{r}\qquad \textbf{(D)}\ \frac{10r}{a}\qquad \textbf{(E)}\ \frac{10a}{r}$

 \textbf{Answer: }E



\textbf{Solution 1}

Use dimensional analysis.
 \[\frac{a/6 \text{ feet}}{r \text{ seconds}} \cdot \frac{1 \text{ yard}}{3 \text{ feet}} \cdot \frac{60 \text{ seconds}}{1 \text{ minute}} \cdot 3 \text{ minutes}\]
 \[\frac{10a}{r} \text{ yards}\]The answer is $\boxed{\textbf{(E)}}$.

 



\textbf{Solution 2}

<i>By Ajeet Dubey | \href{https://ioqm.in}{ioqm.in}</i>

 \textbf{Step 1: Find the speed.}

 The car covers $a/6$ feet in $r$ seconds, so its speed is$\frac{a}{6r}\text{ feet per second}.$

 \textbf{Step 2: Convert the three-minute distance to yards.}

 Since three minutes is $180$ seconds, the distance is$\frac{a}{6r}\cdot180=\frac{30a}{r}\text{ feet}=\frac{10a}{r}\text{ yards}.$

 Thus the answer is $\boxed{\textbf{(E)}}$.

 


	\item (\textit{AHSME 1961}) If the graphs of $2y+x+3=0$ and $3y+ax+2=0$ are to meet at right angles, the value of $a$ is:


 $\textbf{(A)}\ \pm \frac{2}{3} \qquad \textbf{(B)}\ -\frac{2}{3}\qquad \textbf{(C)}\ -\frac{3}{2} \qquad \textbf{(D)}\ 6\qquad \textbf{(E)}\ -6$

 \textbf{Answer: }E



\textbf{Solution 1}

The slope of the first graph is -1/2. The slope of the second is 2, since it is perpendicular, and it is also -a/3 by rearranging. Thus $a=-6$.  Answer is $\boxed{\textbf{(E)}}$.

 



\textbf{Solution 2}

<i>By Ajeet Dubey | \href{https://ioqm.in}{ioqm.in}</i>

 \textbf{Step 1: Determine the two slopes.}

 Writing the lines in slope-intercept form gives slopes$m_1=-\frac12\qquad\text{and}\qquad m_2=-\frac a3.$

 \textbf{Step 2: Use perpendicularity.}

 Perpendicular nonvertical lines satisfy $m_1m_2=-1$. Thus$\left(-\frac12\right)\left(-\frac a3\right)=-1, \qquad \frac a6=-1, \qquad a=-6.$

 Thus the answer is $\boxed{\textbf{(E)}}$.

 


	\item (\textit{AHSME 1961}) Let the set consisting of the squares of the positive integers be called $u$; thus $u$ is the set $1, 4, 9, 16 \ldots$. 
If a certain operation on one or more members of the set always yields a member of the set, 
we say that the set is closed under that operation. Then $u$ is closed under:


 $\textbf{(A)}\ \text{Addition}\qquad \textbf{(B)}\ \text{Multiplication} \qquad \textbf{(C)}\ \text{Division} \qquad\\ \textbf{(D)}\ \text{Extraction of a positive integral root}\qquad \textbf{(E)}\ \text{None of these}$

 \textbf{Answer: }B



\textbf{Solution 1}

Consider each option, case by case.

 For option A, note that $1 + 4 = 5$, so the set is not closed under addition because $5$ is not a perfect square.

 For option B, note that $9 \cdot 4 = 36$ and $4 \cdot 25 = 100$.  Letting one member of the set be $a^2$ and another member be $b^2$ (where $a$ and $b$ are positive integers), the product of the two members is $(ab)^2$.  Since $ab$ is an integer and $(ab)^2$ is a perfect square, the set is closed under multiplication.

 For option C, note that $9 \div 4 = 2.25$, so the set is not closed under division because $2.25$ is not an integer.

 For option D, note that $\sqrt{4} = 2$, so the set is not closed under extraction of a positive integral root because $2$ is not a perfect square.

 Thus, the answer is $\boxed{\textbf{(B)}}$.

 



\textbf{Solution 2}

<i>By Ajeet Dubey | \href{https://ioqm.in}{ioqm.in}</i>

 Let the set consist of the squares of positive integers.

 \textbf{Multiplication works.}

 If $m^2$ and $n^2$ are in the set, then$m^2n^2=(mn)^2,$which is again the square of a positive integer.

 \textbf{The other operations do not always work.}

 For example, $1+4=5$ is not a square, $1/4$ is not the square of a positive integer, and $\sqrt4=2$ is not such a square. Hence multiplication is the only operation under which the set is closed.

 Thus the answer is $\boxed{\textbf{(B)}}$.

 


	\item (\textit{AHSME 1961}) Let $S=(x-1)^4+4(x-1)^3+6(x-1)^2+4(x-1)+1$. Then $S$ equals:


 $\textbf{(A)}\ (x-2)^4 \qquad \textbf{(B)}\ (x-1)^4 \qquad \textbf{(C)}\ x^4 \qquad \textbf{(D)}\ (x+1)^4 \qquad \textbf{(E)}\ x^4+1$

 \textbf{Answer: }C



\textbf{Solution 1}

Let $y = x-1$.  Substitution results in
 \[S = y^4 + 4y^3 + 6y^2 + 4y + 1\]
 \[S = (y+1)^4\]Substituting back results in
 \[S = x^4\]The answer is $\boxed{\textbf{(C)}}$.  This problem can also be solved by traditionally expanding and combining like terms (though it would take much longer).

 



\textbf{Solution 2}

<i>By Ajeet Dubey | \href{https://ioqm.in}{ioqm.in}</i>

 The coefficients $1,4,6,4,1$ are the binomial coefficients in the expansion of $(u+1)^4$. Taking $u=x-1$, the given expression is$(x-1)^4+4(x-1)^3+6(x-1)^2+4(x-1)+1=((x-1)+1)^4=x^4.$

 Thus the answer is $\boxed{\textbf{(C)}}$.

 


	\item (\textit{AHSME 1961}) When simplified, $\log{8} \div \log{\frac{1}{8}}$ becomes:


 $\textbf{(A)}\ 6\log{2} \qquad \textbf{(B)}\ \log{2} \qquad \textbf{(C)}\ 1 \qquad \textbf{(D)}\ 0\qquad \textbf{(E)}\ -1$

 \textbf{Answer: }E



\textbf{Solution 1}

First, note that $\frac{1}{8} = 8^{-1}$.  That means the expression can be rewritten as
 \[\log{8} \div \log{8^{-1}}\]
 \[\log{8} \cdot \frac{1}{\log{8^{-1}}}\]
 \[\log{8} \cdot \frac{1}{-1 \cdot \log{8}}\]This simplifies to $-1$, which is answer choice $\boxed{\textbf{(E)}}$.

 



\textbf{Solution 2}

$\frac{\log 8}{\log \frac{1}{8}} = \log_{\frac{1}{8}}{8} = \boxed{\textbf{(E)}\ {-1}}$.

 -j314andrews

 



\textbf{Solution 3}

<i>By Ajeet Dubey | \href{https://ioqm.in}{ioqm.in}</i>

 Use the logarithm rule $\log(u^{-1})=-\log u$. Since $1/8=8^{-1}$,$\log\left(\frac18\right)=\log(8^{-1})=-\log8.$Therefore$\frac{\log8}{\log(1/8)}=\frac{\log8}{-\log8}=-1.$

 Thus the answer is $\boxed{\textbf{(E)}}$.

 


	\item (\textit{AHSME 1961}) When simplified, the third term in the expansion of $(\frac{a}{\sqrt{x}}-\frac{\sqrt{x}}{a^2})^6$ is:


 $\textbf{(A)}\ \frac{15}{x}\qquad \textbf{(B)}\ -\frac{15}{x}\qquad \textbf{(C)}\ -\frac{6x^2}{a^9} \qquad \textbf{(D)}\ \frac{20}{a^3}\qquad \textbf{(E)}\ -\frac{20}{a^3}$

 \textbf{Answer: }A



\textbf{Solution 1}

By the \href{https://artofproblemsolving.com/wiki/index.php?title=Binomial\_Theorem}{Binomial Theorem}, the third term in the expansion is
 \[\binom{6}{2}(\frac{a}{\sqrt{x}})^{4}(\frac{-\sqrt{x}}{a^2})^2\]
 \[15 \cdot \frac{a^4}{x^2} \cdot \frac{x}{a^4}\]
 \[\frac{15}{x}\]The answer is $\boxed{\textbf{(A)}}$.

 



\textbf{Solution 2}

<i>By Ajeet Dubey | \href{https://ioqm.in}{ioqm.in}</i>

 Let $u=a/\sqrt{x}$ and $v=-\sqrt{x}/a^2$. In the expansion of $(u+v)^6$, the third term is$\binom62u^4v^2.$Now$u^4=\frac{a^4}{x^2},\qquad v^2=\frac{x}{a^4},$so the third term equals$15\cdot\frac{a^4}{x^2}\cdot\frac{x}{a^4}=\frac{15}{x}.$

 Thus the answer is $\boxed{\textbf{(A)}}$.

 


	\item (\textit{AHSME 1961}) Let the two base angles of a triangle be $A$ and $B$, with $B$ larger than $A$. 
The altitude to the base divides the vertex angle $C$ into two parts, $C_1$ and $C_2$, with $C_2$ adjacent to side $a$. Then:


 $\textbf{(A)}\ C_1+C_2=A+B \qquad \textbf{(B)}\ C_1-C_2=B-A \qquad\\ \textbf{(C)}\ C_1-C_2=A-B \qquad \textbf{(D)}\ C_1+C_2=B-A\qquad \textbf{(E)}\ C_1-C_2=A+B$

 \textbf{Answer: }B



\textbf{Solution 1}


                    \begin{figure}[H]
                        \centering
                        \includegraphics[width=0.5\textwidth]{Solution/AHSME_I_1961_8_Solution_Figure_1.png}
                    \end{figure}
                    

 Noting that side $a$ is opposite of angle $A$, label the diagram as shown.

 From the two right triangles,
 \[B + C_2 = 90\]
 \[A + C_1 = 90\]Substitute to get
 \[B + C_2 = A + C_1\]
 \[B - A = C_1 - C_2\]The answer is $\boxed{\textbf{(B)}}$.

 



\textbf{Solution 2}

<i>By Ajeet Dubey | \href{https://ioqm.in}{ioqm.in}</i>

 The altitude creates two right triangles. In the triangle containing angle $A$,$C_1=90^\circ-A,$and in the triangle containing angle $B$,$C_2=90^\circ-B.$Therefore$C_1-C_2=(90^\circ-A)-(90^\circ-B)=B-A.$

 Thus the answer is $\boxed{\textbf{(B)}}$.

 


	\item (\textit{AHSME 1961}) Let $r$ be the result of doubling both the base and exponent of $a^b$, and $b$ does not equal to $0$. 
If $r$ equals the product of $a^b$ by $x^b$, then $x$ equals:


 $\textbf{(A)}\ a \qquad \textbf{(B)}\ 2a \qquad \textbf{(C)}\ 4a \qquad \textbf{(D)}\ 2\qquad \textbf{(E)}\ 4$

 \textbf{Answer: }C



\textbf{Solution 1}

From the problem, $r = (2a)^{2b}$, so
 \[(2a)^{2b} = a^b \cdot x^b\]
 \[(4a^2)^b = (ax)^b\]
 \[4a^2 = ax\]
 \[x = 4a\]Thus, the answer is $\boxed{\textbf{(C)}}$.

 



\textbf{Solution 2}

<i>By Ajeet Dubey | \href{https://ioqm.in}{ioqm.in}</i>

 Doubling both the base and exponent of $a^b$ gives$r=(2a)^{2b}=(4a^2)^b.$But the problem also writes $r=a^bx^b=(ax)^b$. Hence$(4a^2)^b=(ax)^b.$For the positive quantities involved, the bases are equal, so$4a^2=ax\qquad\Longrightarrow\qquad x=4a.$

 Thus the answer is $\boxed{\textbf{(C)}}$.

 


	\item (\textit{AHSME 1961}) Each side of $\triangle ABC$ is $12$ units. $D$ is the foot of the perpendicular dropped from $A$ on $BC$, 
and $E$ is the midpoint of $AD$. The length of $BE$, in the same unit, is:


 $\textbf{(A)}\ \sqrt{18} \qquad \textbf{(B)}\ \sqrt{28} \qquad \textbf{(C)}\ 6 \qquad \textbf{(D)}\ \sqrt{63} \qquad \textbf{(E)}\ \sqrt{98}$

 \textbf{Answer: }D



\textbf{Solution 1}


                    \begin{figure}[H]
                        \centering
                        \includegraphics[width=0.5\textwidth]{Solution/AHSME_I_1961_10_Solution_Figure_1.png}
                    \end{figure}
                    Note that $\triangle ABC$ is an \href{https://artofproblemsolving.com/wiki/index.php?title=Equilateral\_triangle}{equilateral triangle}.  From the \href{https://artofproblemsolving.com/wiki/index.php?title=Pythagorean\_Theorem}{Pythagorean Theorem} (or by using 30-60-90 triangles), $AD = 6\sqrt{3}$.  That means $DE = 3\sqrt{3}$.  Using the Pythagorean Theorem again, $BE = \sqrt{63}$, which is answer choice $\boxed{\textbf{(D)}}$.

 



\textbf{Solution 2}

<i>By Ajeet Dubey | \href{https://ioqm.in}{ioqm.in}</i>

 In an equilateral triangle of side $12$, the altitude is $6\sqrt3$ and meets the opposite side at its midpoint, so $DB=6$. Since $E$ is the midpoint of the altitude,$DE=3\sqrt3.$Triangle $BDE$ is right, and therefore$BE^2=DB^2+DE^2=6^2+(3\sqrt3)^2=36+27=63.$Thus $BE=\sqrt{63}$.

 Thus the answer is $\boxed{\textbf{(D)}}$.

 


	\item (\textit{AHSME 1961}) Two tangents are drawn to a circle from an exterior point $A$; they touch the circle at points $B$ and $C$ respectively. 
A third tangent intersects segment $AB$ in $P$ and $AC$ in $R$, and touches the circle at $Q$. If $AB=20$, then the perimeter of triangle $APR$ is:


 $\textbf{(A)}\ 42\qquad \textbf{(B)}\ 40.5 \qquad \textbf{(C)}\ 40\qquad \textbf{(D)}\ 39\frac{7}{8} \qquad \textbf{(E)}\ \text{not determined by the given information}$

 \textbf{Answer: }C



\textbf{Solution 1}


                    \begin{figure}[H]
                        \centering
                        \includegraphics[width=0.5\textwidth]{Solution/AHSME_I_1961_11_Solution_Figure_1.png}
                    \end{figure}
                    

 Draw the diagram as shown.  Note that the two tangent lines from a single outside point of a circle have the exact same length, so $AB = AC = 20$, $BP = PQ$, and $QR = CR$.

 The perimeter of the triangle is $AP + PR + AR$.  Note that $PQ + QR = PR$, so from substitution, the perimeter is
 \[AP + PQ + QR + AR\]
 \[AP + BP + CR + AR\]
 \[AB + AC\]Thus, the perimeter of the triangle is $40$, so the answer is $\boxed{\textbf{(C)}}$.

 



\textbf{Solution 2 (Non-rigorous)}

Since $Q$ can be anywhere on the circle between $A$ and $B$, it can basically be "on top" of $A$. Then $R$ will be at the same point as $A$, so $APR$ form a degenerate triable with side length $AB=20$. So its perimeter will be $40$. Since $BP=PQ$ and $QR=CR$ by power of a point, as $AP$ and $AR$ decrease in length, $PR=PQ+QR$ will "grow" to compensate, so the perimeter will stay constant with a value of $\boxed{\textbf{(C)}\ 40}$.

 We can also skip verifying that the perimeter will stay constant, since it seems unlikely that MAA would create a question with $\text{not determined by the given information}$ as the answer.

 



\textbf{Solution 3}

<i>By Ajeet Dubey | \href{https://ioqm.in}{ioqm.in}</i>

 Let the third tangent meet $AB$ and $AC$ at $P$ and $R$, and touch the circle at $Q$. Tangent segments from the same external point are equal, so$AB=AC=20,\qquad PB=PQ,\qquad RC=RQ.$Hence$AP=20-PQ, \qquad AR=20-RQ, \qquad PR=PQ+RQ.$The perimeter of $\triangle APR$ is therefore$AP+AR+PR=(20-PQ)+(20-RQ)+(PQ+RQ)=40.$

 Thus the answer is $\boxed{\textbf{(C)}}$.

 


	\item (\textit{AHSME 1961}) The first three terms of a geometric progression are $\sqrt{2}, \sqrt[3]{2}, \sqrt[6]{2}$. Find the fourth term.


 $\textbf{(A)}\ 1\qquad \textbf{(B)}\ \sqrt[7]{2}\qquad \textbf{(C)}\ \sqrt[8]{2}\qquad \textbf{(D)}\ \sqrt[9]{2}\qquad \textbf{(E)}\ \sqrt[10]{2}$

 \textbf{Answer: }A



\textbf{Solution 1}

After rewriting the radicals as fractional exponents, the sequence is $2^{1/2}, 2^{1/3}, 2^{1/6}$.

 The common ratio of the \href{https://artofproblemsolving.com/wiki/index.php?title=Geometric\_sequence}{geometric sequence} is $\frac{2^{1/3}}{2^{1/2}} = 2^{-1/6}$.  Multiplying that by the third term results in $2^0$.  It simplifies to $1$, so the answer is $\boxed{\textbf{(A)}}$.

 



\textbf{Solution 2}

<i>By Ajeet Dubey | \href{https://ioqm.in}{ioqm.in}</i>

 Write the terms as powers of $2$:$2^{1/2},\quad 2^{1/3},\quad 2^{1/6}.$Their exponents decrease by $1/6$, so the common ratio is $2^{-1/6}$. The next exponent is$\frac16-\frac16=0,$and the next term is $2^0=1$.

 Thus the answer is $\boxed{\textbf{(A)}}$.

 


	\item (\textit{AHSME 1961}) The symbol $|a|$ means $a$ if $a$ is a positive number or zero, and $-a$ if $a$ is a negative number. 
For all real values of $t$ the expression $\sqrt{t^4+t^2}$ is equal to:


 $\textbf{(A)}\ t^3\qquad \textbf{(B)}\ t^2+t\qquad \textbf{(C)}\ |t^2+t|\qquad \textbf{(D)}\ t\sqrt{t^2+1}\qquad \textbf{(E)}\ |t|\sqrt{1+t^2}$

 \textbf{Answer: }E



\textbf{Solution 1}

Factor out the $t^2$ inside the square root.
 \[\sqrt{t^2 \cdot (t^2 + 1)}\]
 \[\sqrt{t^2} \cdot \sqrt{t^2 + 1}\]Remember that $\sqrt{t^2} = |t|$ because square rooting a nonnegative real number will always result in a \href{https://artofproblemsolving.com/wiki/index.php?title=Nonnegative\_number}{nonnegative number}.
 \[|t|\sqrt{t^2 + 1}\]The answer is $\boxed{\textbf{(E)}}$.

 



\textbf{Solution 2}

<i>By Ajeet Dubey | \href{https://ioqm.in}{ioqm.in}</i>

 Factor the expression under the radical:$\sqrt{t^4+t^2}=\sqrt{t^2(t^2+1)}.$Because $\sqrt{t^2}=|t|$ for every real $t$, this becomes$|t|\sqrt{t^2+1}.$The absolute value is essential because a square root is nonnegative.

 Thus the answer is $\boxed{\textbf{(E)}}$.

 


	\item (\textit{AHSME 1961}) A rhombus is given with one diagonal twice the length of the other diagonal. 
Express the side of the rhombus in terms of $K$, where $K$ is the area of the rhombus in square inches.


 $\textbf{(A)}\ \sqrt{K}\qquad \textbf{(B)}\ \frac{1}{2}\sqrt{2K}\qquad \textbf{(C)}\ \frac{1}{3}\sqrt{3K}\qquad \textbf{(D)}\ \frac{1}{4}\sqrt{4K}\qquad \textbf{(E)}\ \text{None of these are correct}$

 \textbf{Answer: }E



\textbf{Solution 1}


                    \begin{figure}[H]
                        \centering
                        \includegraphics[width=0.5\textwidth]{Solution/AHSME_I_1961_14_Solution_Figure_1.png}
                    \end{figure}
                    

 Let one of the diagonals of the \href{https://artofproblemsolving.com/wiki/index.php?title=Rhombus}{rhombus} be $a$ units long, so the other diagonal is $2a$ units long.  That means the area of the rhombus is $\frac{1}{2} \cdot 2a \cdot a = a^2$, so $K = a^2$.

 From the \href{https://artofproblemsolving.com/wiki/index.php?title=Pythagorean\_Theorem}{Pythagorean Theorem}, one of the side lengths of the rhombus is $\sqrt{(\frac{a}{2})^2 + a^2} = \sqrt{\frac{5a^2}{4}} = \frac{\sqrt{5K}}{2}$, so the answer is $\boxed{\textbf{(E)}}$. 

 



\textbf{Solution 2}

<i>By Ajeet Dubey | \href{https://ioqm.in}{ioqm.in}</i>

 Let the diagonals of the rhombus be $d$ and $2d$. Its area is$K=\frac12(d)(2d)=d^2.$The diagonals bisect each other at right angles, so a side of the rhombus is the hypotenuse of a right triangle with legs $d/2$ and $d$. Thus$s=\sqrt{\left(\frac d2\right)^2+d^2}=\frac{d\sqrt5}{2}=\frac{\sqrt{5K}}{2}.$This is not among choices (A) through (D).

 Thus the answer is $\boxed{\textbf{(E)}}$.

 


	\item (\textit{AHSME 1961}) If $x$ men working $x$ hours a day for each of $x$ days produce $x$ articles, then the number of articles 
(not necessarily an integer) produced by $y$ men working $y$ hours a day for each of $y$ days is:


 $\textbf{(A)}\ \frac{x^3}{y^2}\qquad \textbf{(B)}\ \frac{y^3}{x^2}\qquad \textbf{(C)}\ \frac{x^2}{y^3}\qquad \textbf{(D)}\ \frac{y^2}{x^3}\qquad \textbf{(E)}\ y$

 \textbf{Answer: }B



\textbf{Solution 1}

Let $k$ be the number of articles produced per hour per person.  By using dimensional analysis,
 \[\frac{x \text{ hours}}{\text{day}} \cdot x \text{ days} \cdot \frac{k \text{ articles}}{\text{hours} \cdot \text{person}} \cdot x \text{ people} = x \text{ articles}\]Solving this yields $k = \frac{1}{x^2}$.  Using dimensional analysis again, the number of articles produced by $y$ men working $y$ hours a day for $y$ days is
 \[\frac{y \text{ hours}}{\text{day}} \cdot y \text{ days} \cdot \frac{\frac{1}{x^2} \text{ articles}}{\text{hours} \cdot \text{person}} \cdot y \text{ people} = \frac{y^3}{x^2} \text{ articles}\]The answer is $\boxed{\textbf{(B)}}$.

 



\textbf{Solution 2 (Simple logic)}

The question is based on the assumption that each person, each hour, each day, will be produce a constant number of items (maybe fractional).

 So it takes $x$ men $x$ hours to produce $\frac{x}{x}=1$ item in a day.

 In a similar manner, 1 man, 1 hour, for a day, can produce $\frac{1}{x^2}$ items. So $y$ men, $y$ hours a day, for $y$ days produce $\frac{y^3}{x^2}$ items. Therefore, the answer is $\boxed{B}$.

 Dimensional analysis is definitely the most rigid, but if you know the ending units (e.g. you know that density is measured in $g/cm^2$ or something like that, you can just treat is as simple proportions and equations.

 



\textbf{Solution 3}

<i>By Ajeet Dubey | \href{https://ioqm.in}{ioqm.in}</i>

 The first situation uses$x\cdot x\cdot x=x^3$man-hours and produces $x$ articles. Therefore the production rate is$\frac{x}{x^3}=\frac1{x^2}$article per man-hour.

 In the second situation there are $y\cdot y\cdot y=y^3$ man-hours. Hence the number of articles produced is$y^3\cdot\frac1{x^2}=\frac{y^3}{x^2}.$

 Thus the answer is $\boxed{\textbf{(B)}}$.

 \textbf{What's happening here? Why isn't the answer "B" ?} Notice that if we change the problem to $x$ men produce $x$ items a day, $y$ men produces how many items a day, then the answer would be $y$. In this case, it would be a direct variation. However, notice that direct variations only have two factors --- an independent and dependent variable each (cause-effect, $x$-$y$). However, there are 3 factors, not 1, that are contributing to how many items are produced in the original problem. This is a joint variation problem, not a direct variation problem. This is the reason why the answer is $\boxed{B}$ (also see that the base unit (1 man/1 hour/1 day) is $\frac{1}{x^2}$).

 Hope that solves your confusions.

 


	\item (\textit{AHSME 1961}) An altitude $h$ of a triangle is increased by a length $m$. How much must be taken from the corresponding base $b$ 
so that the area of the new triangle is one-half that of the original triangle? 


 $\textbf{(A)}\ \frac{bm}{h+m}\qquad \textbf{(B)}\ \frac{bh}{2h+2m}\qquad \textbf{(C)}\ \frac{b(2m+h)}{m+h}\qquad \textbf{(D)}\ \frac{b(m+h)}{2m+h}\qquad \textbf{(E)}\ \frac{b(2m+h)}{2(h+m)}$

 \textbf{Answer: }E



\textbf{Solution 1}

Let $x$ be the value reduced from the base.  Using the variables from the problem, the equation to solve is
 \[\frac{1}{2} (h + m)(b - x) = \frac{1}{2} \cdot \frac{bh}{2}\]Dividing both sides by $\frac{1}{2}$ and $h+m$ results in
 \[b - x = \frac{bh}{2(h+m)}\]
 \[-x = \frac{bh}{2(h+m)} - b\]
 \[x = b - \frac{bh}{2(h+m)}\]The right side can be written as a single fraction by giving $b$ a common denominator.
 \[x = \frac{2bh + 2bm}{2(h+m)} - \frac{bh}{2(h+m)}\]
 \[x = \frac{bh + 2bm}{2(h+m)}\]
 \[x = \frac{b(2m+h)}{2(h+m)}\]The answer is $\boxed{\textbf{(E)}}$.

 



\textbf{Solution 2}

<i>By Ajeet Dubey | \href{https://ioqm.in}{ioqm.in}</i>

 The original area is $bh/2$. After the altitude is increased by $m$ and the base is reduced by $d$, the new area is$\frac12(b-d)(h+m).$This is half the original area, so$\frac12(b-d)(h+m)=\frac14bh.$Thus$2(b-d)(h+m)=bh,$which gives$d=b-\frac{bh}{2(h+m)}=\frac{b(h+2m)}{2(h+m)}.$

 Thus the answer is $\boxed{\textbf{(E)}}$.

 


	\item (\textit{AHSME 1961}) In the base ten number system the number $526$ means $5 \times 10^2+2 \times 10 + 6$. 
In the Land of Mathesis, however, numbers are written in the base $r$. 
Jones purchases an automobile there for $440$ monetary units (abbreviated m.u). 
He gives the salesman a $1000$ m.u bill, and receives, in change, $340$ m.u. The base $r$ is:


 $\textbf{(A)}\ 2\qquad \textbf{(B)}\ 5\qquad \textbf{(C)}\ 7\qquad \textbf{(D)}\ 8\qquad \textbf{(E)}\ 12$

 \textbf{Answer: }D



\textbf{Solution 1}

If Jones received $340$ m.u. change after paying $1000$ m.u. for something that costs $440$ m.u., then
 \[440_r + 340_r = 1000_r\]This equation can be rewritten as
 \[4r^2 + 4r + 3r^2 + 4r = r^3\]Bring all of the terms to one side to get
 \[r^3 - 7r^2 - 8r = 0\]Factor to get
 \[r(r-8)(r+1)=0\]Since \href{https://artofproblemsolving.com/wiki/index.php?title=Base\_numbers}{base numbers} are always positive, base $r$ is $8$, which is answer choice $\boxed{\textbf{(D)}}$.

 



\textbf{Solution 2}

<i>By Ajeet Dubey | \href{https://ioqm.in}{ioqm.in}</i>

 Let the base be $r$. Converting the two numbers to base ten gives$440_r=4r^2+4r, \qquad 340_r=3r^2+4r,$and$1000_r=r^3.$Therefore$7r^2+8r=r^3.$Since $r>0$, divide by $r$:$r^2-7r-8=0=(r-8)(r+1).$A base must be positive, so $r=8$.

 Thus the answer is $\boxed{\textbf{(D)}}$.

 


	\item (\textit{AHSME 1961}) The yearly changes in the population census of a town for four consecutive years are, 
respectively, 25\% increase, 25\% increase, 25\% decrease, 25\% decrease. 
The net change over the four years, to the nearest percent, is:


 $\textbf{(A)}\ -12 \qquad \textbf{(B)}\ -1 \qquad \textbf{(C)}\ 0 \qquad \textbf{(D)}\ 1\qquad \textbf{(E)}\ 12$

 \textbf{Answer: }A



\textbf{Solution 1}

A 25\% increase means the new population is $\frac{5}{4}$ of the original population.  A 25\% decrease means the new population is $\frac{3}{4}$ of the original population.

 Thus, after four years, the population is $1 \cdot \frac{5}{4} \cdot \frac{5}{4} \cdot \frac{3}{4} \cdot \frac{3}{4} = \frac{225}{256}$ times the original population.

 Thus, the net change is -12\%, so the answer is $\boxed{\textbf{(A)}}$.

 



\textbf{Solution 2}

<i>By Ajeet Dubey | \href{https://ioqm.in}{ioqm.in}</i>

 Two successive increases of $25\%$ multiply the population by $(5/4)^2$, while two successive decreases of $25\%$ multiply it by $(3/4)^2$. The overall multiplier is$\left(\frac54\right)^2\left(\frac34\right)^2=\frac{225}{256}\approx0.879.$Thus the final population is about $87.9\%$ of the original, a decrease of about $12.1\%$.

 Thus the closest choice is $-12\%$, or $\boxed{\textbf{(A)}}$.

 


	\item (\textit{AHSME 1961}) Consider the graphs of $y=2\log{x}$ and $y=\log{2x}$. We may say that:


 $\textbf{(A)}\ \text{They do not intersect}\qquad \\ \textbf{(B)}\ \text{They intersect at 1 point only}\qquad \\ \textbf{(C)}\ \text{They intersect at 2 points only} \qquad \\ \textbf{(D)}\ \text{They intersect at a finite number of points but greater than 2} \qquad \\ \textbf{(E)}\ \text{They coincide}$

 \textbf{Answer: }B



\textbf{Solution 1}

Substitute $y$ into the other equation.
 \[2\log{x} = \log{2x}\]Using a logarithm property,
 \[2\log{x} = \log{2} + \log{x}\]Subtract both sides by $\log{x}$ to get
 \[\log{x}=\log{2}\]Thus, $x=2$.  Since both $2\log{x}$ and $\log{2x}$ are functions, there is only one corresponding y-value, so the answer is $\boxed{\textbf{(B)}}$.

 



\textbf{Solution 2}

<i>By Ajeet Dubey | \href{https://ioqm.in}{ioqm.in}</i>

 A logarithm requires $x>0$. Using $2\log x=\log(x^2)$, the equation becomes$\log(x^2)=\log(2x).$Since the logarithm is one-to-one on positive numbers,$x^2=2x.$Thus $x(x-2)=0$. The value $x=0$ is outside the domain, leaving only $x=2$. Therefore the two graphs meet at exactly one point.

 Thus the answer is $\boxed{\textbf{(B)}}$.

 


	\item (\textit{AHSME 1961}) The set of points satisfying the pair of inequalities $y>2x$ and $y>4-x$ is contained entirely in quadrants:


 $\textbf{(A)}\ \text{I and II}\qquad \textbf{(B)}\ \text{II and III}\qquad \textbf{(C)}\ \text{I and III}\qquad \textbf{(D)}\ \text{III and IV}\qquad \textbf{(E)}\ \text{I and IV}$

 \textbf{Answer: }A



\textbf{Solution 1}


                    \begin{figure}[H]
                        \centering
                        \includegraphics[width=0.5\textwidth]{Solution/AHSME_I_1961_20_Solution_Figure_1.png}
                    \end{figure}
                    Graph both equations.  The solution to the system of inequalities is the magenta area because it is in both of the solutions of each individual inequality.  From the graph, the magenta area is in quadrants $I$ and $II$, so the answer is $\boxed{\textbf{(A)}}$.

 



\textbf{Solution 2}

<i>By Ajeet Dubey | \href{https://ioqm.in}{ioqm.in}</i>

 The inequalities are$y>2x\qquad\text{and}\qquad y>4-x,$so every solution lies above both lines. The boundary lines meet where$2x=4-x,$which gives $x=4/3$ and $y=8/3$. This intersection is the lowest point of the upper boundary, and it is above the $x$-axis. Hence every point satisfying both inequalities has $y>0$.

 Such points can have either positive or negative $x$, so the solution set lies only in Quadrants I and II.

 Thus the answer is $\boxed{\textbf{(A)}}$.

 


	\item (\textit{AHSME 1961}) Medians $AD$ and $CE$ of $\triangle ABC$ intersect in $M$. The midpoint of $AE$ is $N$. 
Let the area of $\triangle MNE$ be $k$ times the area of $\triangle ABC$. Then $k$ equals:


 $\textbf{(A)}\ \frac{1}{6}\qquad \textbf{(B)}\ \frac{1}{8}\qquad \textbf{(C)}\ \frac{1}{9}\qquad \textbf{(D)}\ \frac{1}{12}\qquad \textbf{(E)}\ \frac{1}{16}$

 \textbf{Answer: }D



\textbf{Solution}


                    \begin{figure}[H]
                        \centering
                        \includegraphics[width=0.5\textwidth]{Solution/AHSME_I_1961_21_Solution_Figure_1.png}
                    \end{figure}
                    

 Let $[NME] = a$ and $[BMD] = b$.  The altitudes and bases of $\triangle ANM$ and $\triangle NME$ are the same, so $[NME] = [ANM] = a$.  Since the base of $\triangle BEM$ is twice as long as $\triangle NME$ and the altitudes of both triangles are the same, $[BEM] = 2a$.

 Since the altitudes and bases of $\triangle BMD$ and $\triangle CMD$ are the same, $[BMD] = [CMD] = b$.  Additionally, since $\triangle AME$ and $\triangle EMB$ have the same area, $\triangle AMC$ and $\triangle BMC$ also have the same area, so $[AMC] = 2b$.

 Also, the altitudes and bases of $\triangle ABD$ and $\triangle ADC$ are the same, so the area of the two triangles are the same.  Thus,
 \[a+a+2a+b=2b+b\]
 \[4a=2b\]
 \[2a = b\]Thus, the area of $\triangle ABC$ is $12a$, so the area of $\triangle NME$ is $\frac{1}{12}$ the area of $\triangle ABC$.  The answer is $\boxed{\textbf{(D)}}$.

 


 


	\item (\textit{AHSME 1961}) If $3x^3-9x^2+kx-12$ is divisible by $x-3$, then it is also divisible by:


 $\textbf{(A)}\ 3x^2-x+4\qquad \textbf{(B)}\ 3x^2-4\qquad \textbf{(C)}\ 3x^2+4\qquad \textbf{(D)}\ 3x-4 \qquad \textbf{(E)}\ 3x+4$

 \textbf{Answer: }C



\textbf{Solution}

If $3x^3-9x^2+kx-12$ is divisible by $x-3$, then by the \href{https://artofproblemsolving.com/wiki/index.php?title=Remainder\_Theorem}{Remainder Theorem}, plugging in $3$ in the cubic results in $0$.
 \[3 \cdot 3^3 - 9 \cdot 3^2 + 3k - 12 = 0\]Combine like terms to get
 \[3k - 12 = 0\]Thus, $k=4$.  The cubic is $3x^3-9x^2+4x-12$, and it can be factored (by grouping or synthetic division) into
 \[(3x^2+4)(x-3)\]Thus, the answer is $\boxed{\textbf{(C)}}$.

 


	\item (\textit{AHSME 1961}) Points $P$ and $Q$ are both in the line segment $AB$ and on the same side of its midpoint. $P$ divides $AB$ in the ratio $2:3$, 
and $Q$ divides $AB$ in the ratio $3:4$. If $PQ=2$, then the length of $AB$ is:


 $\textbf{(A)}\ 60\qquad \textbf{(B)}\ 70\qquad \textbf{(C)}\ 75\qquad \textbf{(D)}\ 80\qquad \textbf{(E)}\ 85$

 \textbf{Answer: }B



\textbf{Solution}


                    \begin{figure}[H]
                        \centering
                        \includegraphics[width=0.5\textwidth]{Solution/AHSME_I_1961_23_Solution_Figure_1.png}
                    \end{figure}
                    

 Draw diagram as shown, where $P$ and $Q$ are on the same side.  Let $AP = x$ and $QB = y$.

 Since $P$ divides $AB$ in the ratio $2:3$, $\frac{x}{y+2} = \frac{2}{3}$.  Since $Q$ divides $AB$ in the ratio $3:4$, $\frac{x+2}{y} = \frac{3}{4}$.  Cross multiply to get a \href{https://artofproblemsolving.com/wiki/index.php?title=System\_of\_equations}{system of equations}.
 \[3x = 2y+4\]
 \[4x+8=3y\]Solve the system to get $x = 28$ and $y = 40$.  Thus, $AB = 40 + 28 + 2 = 70$, so the answer is $\boxed{\textbf{(B)}}$.

 


	\item (\textit{AHSME 1961}) Thirty-one books are arranged from left to right in order of increasing prices. 
The price of each book differs by $\textdollar 2$ from that of each adjacent book. 
For the price of the book at the extreme right a customer can buy the middle book and the adjacent one. Then:


 $\textbf{(A)}\ \text{The adjacent book referred to is at the left of the middle book}\qquad \\ \textbf{(B)}\ \text{The middle book sells for \textdollar 36} \qquad \\ \textbf{(C)}\ \text{The cheapest book sells for \textdollar4} \qquad \\ \textbf{(D)}\ \text{The most expensive book sells for \textdollar64 } \qquad \\ \textbf{(E)}\ \text{None of these is correct }$

 \textbf{Answer: }A



\textbf{Solution}

Let the leftmost book cost $c$ dollars.  The 16th book (the middle book) costs $c+30$ dollars, and the 31st book costs $c+60$ dollars.

 If the book to the left of the middle book is used, the equation to solve for the price of the leftmost book is $c+30+c+29=c+60$.  Solving yields $c=1$, so the price of the middle book is $\textdollar{31}$ and the price of the rightmost book is $\textdollar{61}$.

 If the book to the right of the middle book is used, the equation to solve for the price of the leftmost book is $c+30+c+31=c+60$.  Solving yields $c=-1$, which can not happen.

 Thus, the answer is $\boxed{\textbf{(A)}}$.

 


	\item (\textit{AHSME 1961}) $\triangle ABC$ is isosceles with base $AC$. Points $P$ and $Q$ are respectively in $CB$ and $AB$ and such that $AC=AP=PQ=QB$. 
The number of degrees in $\angle B$ is:


 $\textbf{(A)}\ 25\frac{5}{7}\qquad \textbf{(B)}\ 26\frac{1}{3}\qquad \textbf{(C)}\ 30\qquad \textbf{(D)}\ 40\qquad \textbf{(E)}\ \text{Not determined by the information given}$

 \textbf{Answer: }A



\textbf{Solution}


                    \begin{figure}[H]
                        \centering
                        \includegraphics[width=0.5\textwidth]{Solution/AHSME_I_1961_25_Solution_Figure_1.png}
                    \end{figure}
                    

 Let $b$ be the measure of $\angle B$.  $\triangle BQP$ is an \href{https://artofproblemsolving.com/wiki/index.php?title=Isosceles\_triangle}{isosceles triangle}, so $\angle BPQ = b$ and $\angle BQP = 180 - 2b$.

 $AB$ is a line, so $\angle AQP = 2b$.  Since $\triangle QPA$ is isosceles as well, $\angle QAP = 2b$ and $\angle QPA = 180 - 4b$.

 $BC$ is a line, so $\angle APC = 3b$.  Since $\triangle APC$ is isosceles as well, $\angle ACB = 3b$.  $\triangle ABC$ is also isosceles, so $\angle BAC = 3b$, so $\angle PAC = b$.

 The angles in a triangle add up to $180$ degrees, so $3b + 3b + b = 180$.  Solving the equation yields $b = \frac{180}{7}$.  Thus, $\angle B = 25\frac{5}{7}^{\circ}$, so the answer is $\boxed{\textbf{(A)}}$.

 


	\item (\textit{AHSME 1961}) For a given arithmetic series the sum of the first $50$ terms is $200$, and the sum of the next $50$ terms is $2700$. 
The first term in the series is:


 $\textbf{(A)}\ -1221 \qquad \textbf{(B)}\ -21.5 \qquad \textbf{(C)}\ -20.5 \qquad \textbf{(D)}\ 3 \qquad \textbf{(E)}\ 3.5$

 \textbf{Answer: }C



\textbf{Solution 1}

Let the first term of the \href{https://artofproblemsolving.com/wiki/index.php?title=Arithmetic\_sequence}{arithmetic sequence} be $a$ and the common difference be $d$.

 The $50^{\text{th}}$ term of the sequence is $a+49d$, so the sum of the first $50$ terms is $\frac{50(a + a + 49d)}{2}$.

 The $51^{\text{th}}$ term of the sequence is $a+50d$ and the $100^{\text{th}}$ term of the sequence is $a+99d$, so the sum of the next $50$ terms is $\frac{50(a+50d+a+99d)}{2}$.

 Substituting in values results in this \href{https://artofproblemsolving.com/wiki/index.php?title=System\_of\_equations}{system of equations}.
 \[2a+49d=8\]
 \[2a+149d=108\]Solving for $a$ yields $a = \frac{-41}{2}$.  The first term is $-20.5$, which is answer choice $\boxed{\textbf{(C)}}$.

 



\textbf{Solution 2}

Same as solution 1 but faster.

 Theory:In an AP, we can write$S_n = An^2+Bn$, where common difference($d) = 2A$ and first term($a) = A(1^2)+B(1) = A+B$ 

 This is already known, and not part of the solution :) .

 So, given:

 $S_{50} = 2500A+50B=200$
$S_{100} - S_{50} = 7500A+50B=2700$
$\implies A=0.5 \text{ and } B = -21$
$\implies \text{first term} = A+B = \boxed{\textbf{ -20.5 (C) }}$



	\item (\textit{AHSME 1961}) Given two equiangular polygons $P_1$ and $P_2$ with different numbers of sides; 
each angle of $P_1$ is $x$ degrees and each angle of $P_2$ is $kx$ degrees, 
where $k$ is an integer greater than $1$. 
The number of possibilities for the pair $(x, k)$ is:


 $\textbf{(A)}\ \infty\qquad \textbf{(B)}\ \text{finite, but greater than 2}\qquad \textbf{(C)}\ 2\qquad \textbf{(D)}\ 1\qquad \textbf{(E)}\ 0$

 \textbf{Answer: }D



\textbf{Solution}

Each angle in each equiangular polygon is $\frac{180(n-2)}{n} = 180 - \frac{360}{n}$, where $n$ is the number of sides.  As $n$ gets larger, each angle in the equiangular polygon approaches (but does not reach) $180^{\circ}$.  That means $kx < 180$, so $x < \frac{180}{k}$.

Recall that each angle in an equiangular triangle has $60^{\circ}$, each angle in an equiangular quadrilateral has $90^{\circ}$, and so on.  If $k = 2$, then $x < 90$.  That means the only option available is $x = 60$ because an equiangular triangle is the only equiangular polygon to have acute angles.  If $k = 3$, then $x < 60$, and there are no values of $x$ that satisfy.  As $k$ increases, $\frac{180}{k}$ decreases even further, so we can for sure say that there is only 1 possibility for the pair $(x, k)$, which is answer choice $\boxed{\textbf{(D)}}$.

 


	\item (\textit{AHSME 1961}) If $2137^{753}$ is multiplied out, the units' digit in the final product is:


 $\textbf{(A)}\ 1\qquad \textbf{(B)}\ 3\qquad \textbf{(C)}\ 5\qquad \textbf{(D)}\ 7\qquad \textbf{(E)}\ 9$

 \textbf{Answer: }D



\textbf{Solution}

$7^1$ has a unit digit of $7$.$7^2$ has a unit digit of $9$.$7^3$ has a unit digit of $3$.$7^4$ has a unit digit of $1$.$7^5$ has a unit digit of $7$.

 Notice that the unit digit eventually cycles to itself when the exponent is increased by $4$.  It also does not matter what the other digits are in the base because the units digit is found by multiplying by only the units digit.  Since $753$ leaves a remainder of $1$ after being divided by $4$, the units digit of $2137^{753}$ is $7$, which is answer choice $\boxed{\textbf{(D)}}$.

 \textbf{Alternate Solution} 
            \begin{itemize}
        \item $Lemma$ ($Fermat's$ $Theorem$): If $p$ is a prime and $a$ is an integer prime to $p$ then we have $a^{p-1} \equiv 1\ (\textrm{mod}\ p)$.
            \end{itemize}
            
            \begin{itemize}
        \item Let's define $U$($x$) as units digit funtion of $x$.
            \end{itemize}
            We can clearly observe that,

 \begin{verbatim}
$U$($7^1$)= $7$  .      .  .      .  .      . $U$($7^4$)= $1$
\end{verbatim}

and we can see by Fermat's Theorem that this cycle repeats with the cyclicity of $4$ .                                                     Now $753$ = $4k$ + $1$ $=>$ $U($7$^{753}$)$=$7$.$ ~GEOMETRY-WIZARD

 


	\item (\textit{AHSME 1961}) Let the roots of $ax^2+bx+c=0$ be $r$ and $s$. The equation with roots $ar+b$ and $as+b$ is:


 $\textbf{(A)}\ x^2-bx-ac=0\qquad \textbf{(B)}\ x^2-bx+ac=0 \qquad\\ \textbf{(C)}\ x^2+3bx+ca+2b^2=0 \qquad \textbf{(D)}\ x^2+3bx-ca+2b^2=0 \qquad\\ \textbf{(E)}\ x^2+bx(2-a)+a^2c+b^2(a+1)=0$

 \textbf{Answer: }B



\textbf{Solution}

From \href{https://artofproblemsolving.com/wiki/index.php?title=Vieta\%27s\_Formulas}{Vieta's Formulas}, $r+s=\frac{-b}{a}$ and $rs=\frac{c}{a}$ in the original quadratic.

 The sum of the roots in the new quadratic is
 \[ar + b + as + b\]
 \[a(r+s) + 2b\]
 \[b\]The product of the roots in the new quadratic is
 \[(ar + b)(as + b)\]
 \[a^2rs + arb + asb + b^2\]
 \[a^2rs + ab(r+s) + b^2\]
 \[ac - b^2 + b^2\]
 \[ac\]Thus, the new quadratic is $x^2-bx+ac$.  The answer is $\boxed{\textbf{(B)}}$.

 



\textbf{Solution 2}

Let $f(x)=ax^2+bx+c$ and $g(x)$ be the desired function with roots $ar+b$ and $as+b$. Applying graphical transformations, we have
 \[f(r)=f(s)=0\]
 \[f(\frac{(ar+b)-b}{a})=f(\frac{(as+b)-b}{a})=0\]Thus $ar+b$ and $as+b$ are the roots to
 \[g(x)=f(\frac{x-b}{a})\]Plugging $x=\frac{x-b}{a}$ into $f(x)$ gets $g(x)=x^2-bx+ac \boxed{B}$.

 


	\item (\textit{AHSME 1961}) If $\log_{10}2=a$ and $\log_{10}3=b$, then $\log_{5}12$ equals:


 $\textbf{(A)}\ \frac{a+b}{1+a}\qquad \textbf{(B)}\ \frac{2a+b}{1+a}\qquad \textbf{(C)}\ \frac{a+2b}{1+a}\qquad \textbf{(D)}\ \frac{2a+b}{1-a}\qquad \textbf{(E)}\ \frac{a+2b}{1-a}$

 \textbf{Answer: }D



\textbf{Solution}

From the \href{https://artofproblemsolving.com/wiki/index.php?title=Change\_of\_base\_formula}{change of base formula}, $\log_{5}12 = \frac{\log_{10}12}{\log_{10}5}$.

 For the numerator, $2 \cdot 2 \cdot 3 = 12$, so $\log_{10}12 = \log_{10}2 + \log_{10}2 + \log_{10}3 = 2a+b$.

 For the denominator, note that $\log_{10}10 = 1$.  Thus, $\log_{10}5 = \log_{10}10 - \log_{10}2 = 1 - a$.

 Thus, the fraction is $\frac{2a+b}{1-a}$, so the answer is $\boxed{\textbf{(D)}}$.

 


	\item (\textit{AHSME 1961}) In $\triangle ABC$ the ratio $AC:CB$ is $3:4$. The bisector of the exterior angle at $C$ intersects $BA$ extended at $P$ 
($A$ is between $P$ and $B$). The ratio $PA:AB$ is:


 $\textbf{(A)}\ 1:3 \qquad \textbf{(B)}\ 3:4 \qquad \textbf{(C)}\ 4:3 \qquad \textbf{(D)}\ 3:1 \qquad \textbf{(E)}\ 7:1$

 \textbf{Answer: }D



\textbf{Solution}


                    \begin{figure}[H]
                        \centering
                        \includegraphics[width=0.5\textwidth]{Solution/AHSME_I_1961_31_Solution_Figure_1.png}
                    \end{figure}
                    Let $AC = 3n$ and $BC = 4n$.  Draw $X$, where $X$ is on $BC$ and $AC \parallel PX$.  By AA Similarity, $\triangle ABC \sim \triangle PBX$, so $PX = 3an$, $BX = 4an$, and $CX = 4an - 4n$.

 
Also, let $\angle ABC = a$ and $\angle BAC = b$.  Since the angles of a triangle add up to $180^{\circ}$, $\angle BCA = 180-a-b$.  By Exterior Angle Theorem, $\angle ACX = a+b$, and since $CP$ bisects $\angle ACX$, $\angle PCX = \frac{a+b}{2}$.  Because $AC \parallel PX$, $\angle BXP = 180 - a - b$.  Thus, $\angle CPX = \frac{a+b}{2}$, making $\triangle CPX$ an \href{https://artofproblemsolving.com/wiki/index.php?title=Isosceles\_triangle}{isosceles triangle}.

 
Because $\triangle CPX$ is isosceles, $PX = CX$, so $4an - 4n = 3an$.  That means $a = 4$, so $PB = 4 \cdot AB$.  Thus, $PA = PB - AB = 3 \cdot AB$, so $PA : AB = 3:1$.  The answer is $\boxed{\textbf{(D)}}$, and it can be verified (or obtained) by making $ABC$ a 3-4-5 right triangle.

 


	\item (\textit{AHSME 1961}) A regular polygon of $n$ sides is inscribed in a circle of radius $R$. The area of the polygon is $3R^2$. Then $n$ equals: 


 $\textbf{(A)}\ 8\qquad \textbf{(B)}\ 10\qquad \textbf{(C)}\ 12\qquad \textbf{(D)}\ 15\qquad \textbf{(E)}\ 18$

 \textbf{Answer: }C



\textbf{Solution}

Note that the distance from the center of the \href{https://artofproblemsolving.com/wiki/index.php?title=Circle}{circle} to each of the vertices of the inscribed \href{https://artofproblemsolving.com/wiki/index.php?title=Regular\_polygon}{regular polygon} equals the radius $R$.  Since each side of a regular polygon is the same length, all the angles between the two lines from the center to the two vertices of a side is the same.

 That means each of these angles between the two lines from the center to the two vertices of a side equals $\frac{360}{n}$ degrees.  Thus, the area of the polygon is
 \[n \cdot \frac{1}{2}R^2\sin\left(\frac{360}{n}^{\circ}\right) = 3R^2\]Dividing both sides by $R^2$ yields
 \[\frac{n}{2}\sin\left(\frac{360}{n}^{\circ}\right) = 3\]Multiply both sides by $\frac{2}{n}$ to get
 \[\sin\left(\frac{360}{n}^{\circ}\right) = \frac{6}{n}\]At this point, use trial-and-error for each of the answer choices.  When checking $n = 12$, the equation results in $\sin(30^{\circ}) = \frac{1}{2}$, which is correct.  Thus, the answer is $\boxed{\textbf{(C)}}$.

 


	\item (\textit{AHSME 1961}) The number of solutions of $2^{2x}-3^{2y}=55$, in which $x$ and $y$ are integers, is:


 $\textbf{(A)}\ 0\qquad \textbf{(B)}\ 1\qquad \textbf{(C)}\ 2\qquad \textbf{(D)}\ 3\qquad \textbf{(E)}\ \text{More than three, but finite}$

 \textbf{Answer: }B



\textbf{Solution 1}

Let $a = 2^x$ and $b = 3^y$.  Substituting these values results in
 \[a^2 - b^2 = 55\]Factor the \href{https://artofproblemsolving.com/wiki/index.php?title=Difference\_of\_squares}{difference of squares} to get
 \[(a + b)(a - b) = 55\]If $y < 0$, then $55 + 3^{2y} < 64$, so $y$ can not be negative.  If $x < 0$, then $2^{2x} < 1$.  Since $3^{2y}$ is always positive, the result would be way less than $55$, so $x$ can not be negative.  Thus, $x$ and $y$ have to be \href{https://artofproblemsolving.com/wiki/index.php?title=Nonnegative}{nonnegative}, so $a$ and $b$ are \href{https://artofproblemsolving.com/wiki/index.php?title=Integers}{integers}.  Thus,
 \[a+b=55 \text{ and } a-b=1\]
 \[\text{or}\]
 \[a+b=11 \text{ and } a-b=5\]From the first case, $a = 28$ and $b = 27$.  Since $2^x = 28$ does not have an integral solution, the first case does not work.From the second case, $a = 8$ and $b = 3$.  Thus, $x = 3$ and $y = 1$.  Thus, there is only one solution, which is answer choice $\boxed{\textbf{(B)}}$.

 



\textbf{Solution 2}

<i>By Ajeet Dubey | \href{https://ioqm.in}{ioqm.in}</i>

 \textbf{Step 1: Eliminate negative exponents}

 Suppose $x<0$. Then
 \[0<2^{2x}<1 \qquad\text{and}\qquad 3^{2y}>0.\]Therefore
 \[2^{2x}-3^{2y}<1,\]which cannot equal $55$. Hence $x\geq 0$.

 Since $x\geq 0$, the number $2^{2x}$ is an integer. Thus
 \[3^{2y}=2^{2x}-55\]must also be an integer. However, if $y<0$, then
 \[3^{2y}=\frac{1}{3^{-2y}}\]is a proper fraction. This is impossible, so $y\geq 0$.

 \textbf{Step 2: Factor the equation}

 Set
 \[a=2^x \qquad\text{and}\qquad b=3^y.\]Now $a$ and $b$ are positive integers, and
 \[a^2-b^2=55.\]Factoring gives
 \[(a-b)(a+b)=55.\]

 \textbf{Step 3: Check the factor pairs}

 Because $a>b>0$, both factors are positive. The only possibilities are
 \[(a-b,a+b)=(1,55) \qquad\text{or}\qquad (5,11).\]

            \begin{itemize}
        \item If $(a-b,a+b)=(1,55)$, then $(a,b)=(28,27)$. This is impossible because $28$ is not a power of $2$.
        \item If $(a-b,a+b)=(5,11)$, then $(a,b)=(8,3)$. Therefore
            \end{itemize}
            
 \[2^x=8 \qquad\text{and}\qquad 3^y=3,\]so
 \[x=3 \qquad\text{and}\qquad y=1.\]

 \textbf{Conclusion}

 The only integer solution is
 \[\boxed{(x,y)=(3,1)}.\]Therefore, the number of solutions is $\boxed{1}$, and the answer is $\boxed{\textbf{(B)}}$.

 


	\item (\textit{AHSME 1961}) Let S be the set of values assumed by the function 
 \[\frac{2x+3}{x+2}\]
when $x$ is any member of the interval $x \ge 0$. If there exists a number $M$ such that no number of the set $S$ is greater than $M$, 
then $M$ is an upper bound of $S$. If there exists a number $m$ such that no number of the set $S$ is less than $m$, 
then $m$ is a lower bound of $S$. We may then say:


 $\textbf{(A)}\ \text{m is in S, but M is not in S}\qquad\\ \textbf{(B)}\ \text{M is in S, but m is not in S}\qquad\\ \textbf{(C)}\ \text{Both m and M are in S}\qquad\\ \textbf{(D)}\ \text{Neither m nor M are in S}\qquad\\ \textbf{(E)}\ \text{M does not exist either in or outside S}$

 \textbf{Answer: }A



\textbf{Solution}

This problem is really finding the \href{https://artofproblemsolving.com/wiki/index.php?title=Range}{range} of a function with a restricted \href{https://artofproblemsolving.com/wiki/index.php?title=Domain\_(function)}{domain}.

 


 Dividing $x+2$ into $2x+3$ yields $2 - \frac{1}{x+2}$.  Since $x \ge 0$, as $x$ gets larger, $\frac{1}{x+2}$ approaches $0$, so $\frac{2x+3}{x+2}$ approaches $2$ as $x$ gets larger.  That means $M = 2$.  Since $\frac{1}{x+2}$ can never be $0$, $\frac{2x+3}{x+2}$ can never be $2$, so $M$ is not in the set $S$.

 


 For the smallest value, plug in $0$ in $\frac{2x+3}{x+2}$ to get $\frac{3}{2}$, so $m = \frac{3}{2}$.  Since plugging in $0$ results in $m$, $m$ is in the set $S$.

 


 Thus, the answer is $\boxed{\textbf{(A)}}$.

 


	\item (\textit{AHSME 1961}) The number $695$ is to be written with a factorial base of numeration, that is, $695=a_1+a_2\times2!+a_3\times3!+ \cdots a_n \times n!$ 
where $a_1, a_2, a_3 \cdots a_n$ are integers such that $0 \le a_k \le k$, and $n!$ means $n(n-1)(n-2)\cdots2 \times 1$. Find $a_4$
$\textbf{(A)}\ 0\qquad \textbf{(B)}\ 1\qquad \textbf{(C)}\ 2\qquad \textbf{(D)}\ 3\qquad \textbf{(E)}\ 4$

 \textbf{Answer: }D



\textbf{Solution}

This problem can be approached similarly to other \href{https://artofproblemsolving.com/wiki/index.php?title=Base\_number}{base number} problems.

 Since $120 < 695 < 720$, divide $695$ by $120$.  The quotient is $5$ and the remainder is $95$, so rewrite the number as
 \[695 = 5 \cdot 120 + 95\]Similarly, dividing $95$ by $24$ results in a quotient of $3$ and a remainder of $23$, so the number can be rewritten as
 \[695 = 5 \cdot 120 + 3 \cdot 24 + 23\]Repeat the steps to get
 \[695 = 5 \cdot 120 + 3 \cdot 24 + 3 \cdot 6 + 2 \cdot 2 + 1\]The answer is $\boxed{\textbf{(D)}}$.  One can also stop at the second step by noting $23 < 24$.

 


	\item (\textit{AHSME 1961}) In $\triangle ABC$ the median from $A$ is given perpendicular to the median from $B$. If $BC=7$ and $AC=6$, find the length of $AB$. 


 $\textbf{(A)}\ 4\qquad \textbf{(B)}\ \sqrt{17} \qquad \textbf{(C)}\ 4.25\qquad \textbf{(D)}\ 2\sqrt{5} \qquad \textbf{(E)}\ 4.5$

 \textbf{Answer: }B



\textbf{Solution}


                    \begin{figure}[H]
                        \centering
                        \includegraphics[width=0.5\textwidth]{Solution/AHSME_I_1961_36_Solution_Figure_1.png}
                    \end{figure}
                    By SAS Similarity, $\triangle ABC \sim \triangle MNC$, so $AB \parallel MN$.  Thus, by AA Similarity, $\triangle AGB \sim \triangle NGM$.

 Let $a = GN$ and $b = GM$, so $AG = 2a$ and $BG = 2b$.  By the \href{https://artofproblemsolving.com/wiki/index.php?title=Pythagorean\_Theorem}{Pythagorean Theorem},
 \[4a^2 + b^2 = 9\]
 \[a^2 + 4b^2 = \frac{49}{4}\]Adding the two equations yields $5a^2 + 5b^2 = \frac{85}{4}$, so $a^2 + b^2 = \frac{17}{4}$.  Thus, $MN = \frac{\sqrt{17}}{2}$, so $AB = \sqrt{17}$, which is answer choice $\boxed{\textbf{(B)}}$.

 


	\item (\textit{AHSME 1961}) In racing over a given distance $d$ at uniform speed, $A$ can beat $B$ by $20$ yards, $B$ can beat $C$ by $10$ yards, 
and $A$ can beat $C$ by $28$ yards. Then $d$, in yards, equals:


 $\textbf{(A)}\ \text{Not determined by the given information}\qquad \textbf{(B)}\ 58\qquad \textbf{(C)}\ 100\qquad \textbf{(D)}\ 116\qquad \textbf{(E)}\ 120$

 \textbf{Answer: }C



\textbf{Solution}

Let $a$ be speed of $A$, $b$ be speed of $B$, and $c$ be speed of $C$.

 Person $A$ finished the track in $\frac{d}{a}$ minutes, so $B$ traveled $\frac{db}{a}$ yards at the same time.  Since $B$ is $20$ yards from the finish line, the first equation is
 \[\frac{db}{a} + 20 = d\]Using similar steps, the second and third equation are, respectively,
 \[\frac{dc}{b} + 10 = d\]
 \[\frac{dc}{a} + 28 = d\]Get rid of the denominator in all three equations by multiplying both sides by the denominator in each equation.
 \[db + 20a = da\]
 \[dc + 10b = db\]
 \[dc + 28a = da\]These equations can be rearranged to get the following.
 \[(d-20)a = db\]
 \[(d-10)b = dc\]
 \[(d-28)a = dc\]

 Solving for $b$ in the first equation yields $\frac{(d-20)a}{d} = b$.  Substituting for $b$ and solving for $c$ in the second equation yields $\frac{(d-10)(d-20)a}{d^2} = c$.  Substituting for $c$ in the third equation yields
 \[(d-28)a = \frac{(d-10)(d-20)a}{d}\]
 \[d^2 - 28d = d^2 - 30d + 200\]Solve the equation to get $d = 100$.  Thus, the track is $100$ yards long, which is answer choice $\boxed{\textbf{(C)}}$.

 


 



\textbf{Solution 2}

Let $s_a$, $s_b$, $s_c$ be the speeds of $A$, $B$, $C$, respectively and let $t_a$, $t_b$, $t_c$ be their respective times needed, such that 
 \[d=s_at_a=s_bt_b=s_ct_b\]Putting each pair of them into a race gets the $3$ equations
 \[s_at_a-s_bt_a=20\]
 \[s_bt_b-s_ct_b=10\]
 \[s_at_a-s_ct_a=28\]Subtract the first equation from the third yields
 \[s_bt_a-s_ct_a=t_a(s_b-s_c)=8\]Divide the equation by the second original equation yields 
 \[\frac{t_a}{t_b}=\frac{4}{5}\]Since $s_at_a=s_bt_b$ we have
 \[\frac{s_a}{s_b}=\frac{5}{4}\]Thus from the first equation we have
 \[s_at_a-s_bt_a=s_at_a-\frac{4}{5}s_at_a=20\Longrightarrow d=s_at_a=100 \boxed{(C)}\]



	\item (\textit{AHSME 1961}) $\triangle ABC$ is inscribed in a semicircle of radius $r$ so that its base $AB$ coincides with diameter $AB$. 
Point $C$ does not coincide with either $A$ or $B$. Let $s=AC+BC$. Then, for all permissible positions of $C$:


 $\textbf{(A)}\ s^2\le8r^2\qquad \textbf{(B)}\ s^2=8r^2 \qquad \textbf{(C)}\ s^2 \ge 8r^2 \qquad\\ \textbf{(D)}\ s^2\le4r^2 \qquad \textbf{(E)}\ s^2=4r^2$

 \textbf{Answer: }A



\textbf{Solution 1}


                    \begin{figure}[H]
                        \centering
                        \includegraphics[width=0.5\textwidth]{Solution/AHSME_I_1961_38_Solution_Figure_1.png}
                    \end{figure}
                    Since $s=AC+BC$, $s^2 = AC^2 + 2 \cdot AC \cdot BC + BC^2$.  Since $\triangle ABC$ is inscribed and $AB$ is the diameter, $\triangle ABC$ is a right triangle, and by the \href{https://artofproblemsolving.com/wiki/index.php?title=Pythagorean\_Theorem}{Pythagorean Theorem}, $AC^2 + BC^2 = AC^2 = (2r)^2$.  Thus, $s^2 = 4r^2 + 2 \cdot AC \cdot BC$.

 


 The area of $\triangle ABC$ is $\frac{AC \cdot BC}{2}$, so $2 \cdot [ABC] = AC \cdot BC$.  That means $s^2 = 4r^2 + 4 \cdot [ABC]$.  The area of $\triangle ABC$ can also be calculated by using base $AB$ and the altitude from $C$.  The maximum possible value of the altitude is $r$, so the maximum area of $\triangle ABC$ is $r^2$.

 


 Therefore, $s^2 \le 8r^2$, so the answer is $\boxed{\textbf{(A)}}$.

 



\textbf{Solution 2 (Mean Inequality Chain)}

First $AB=2r$. $AC^2+BC^2=4r^2$.

 By quickly considering extreme cases (when $C$ coincides with $A$ or $B$) we can suspect that the maximum value of $S$ is achieved when $CO$ is perpendicular to $AB$, where $O$ is the origin of the semicircle. We will prove this using the Mean Inequality Chain. Let $AC=x, BC=y$. $x^2+y^2=4r^2$. By QM-AM, $\sqrt{\frac{x^2+y^2}{2}}\ge \frac{x+y}{2}=0.5s$. Plug in $x^2+y^2=4r^2$, we have $2\sqrt{2}r\ge s$. Square both sides, we find that $\boxed{\textbf{(A)}}$ is correct.

 



\textbf{Solution 3 (Calculus Optimization)}

This solution utilizes optimization techniques taught in Calculus 1 classes.

 Let the equation of the semicircle be $x^2+y^2=r^2, y\ge 0$. Notice that in this semicircle, $y$ is a function of $x$.$y=\sqrt{r^2-x^2}$.

 Now let $C(x, \sqrt{r^2-x^2})$. Using the Pythagorean Theorem with $A (-r,0), B(r,0)$, finding $s$ in terms of $x$, differentiating and finding the absolute maximum, we can find that $\boxed{\textbf{(A)}}$ is correct. 

 P.S.: The process of bash is omitted here since anybody that actually learned optimization problems and paid attention to it (this is core curriculum in regular college math) should be able to bash it out.

 


	\item (\textit{AHSME 1961}) Any five points are taken inside or on a square with side length $1$. Let $a$ be the smallest possible number with the 
property that it is always possible to select one pair of points from these five such that the distance between them 
is equal to or less than $a$. Then $a$ is:


 $\textbf{(A)}\ \sqrt{3}/3\qquad \textbf{(B)}\ \sqrt{2}/2\qquad \textbf{(C)}\ 2\sqrt{2}/3\qquad \textbf{(D)}\ 1 \qquad \textbf{(E)}\ \sqrt{2}$

 \textbf{Answer: }B



\textbf{Solution}

Partition the unit square into four smaller squares of sidelength $\frac{1}{2}$. Each of the five points lies in one of these squares, and so by the Pigeonhole Principle, there exists two points in the same $\frac{1}{2}\times \frac{1}{2}$ square - the maximum possible distance between them being $\frac{\sqrt{2}}{2}$ by Pythagoras. 

 Hence our answer is $\fbox{B}$.

 


	\item (\textit{AHSME 1961}) Find the minimum value of $\sqrt{x^2+y^2}$ if $5x+12y=60$. 


 $\textbf{(A)}\ \frac{60}{13}\qquad \textbf{(B)}\ \frac{13}{5}\qquad \textbf{(C)}\ \frac{13}{12}\qquad \textbf{(D)}\ 1\qquad \textbf{(E)}\ 0$

 \textbf{Answer: }A



\textbf{Solution 1}

Let $x^2 + y^2 = r^2$, so $r = \sqrt{x^2 + y^2}$.  Thus, this problem is really finding the shortest distance from the origin to the line $5x + 12y = 60$.

 
                    \begin{figure}[H]
                        \centering
                        \includegraphics[width=0.5\textwidth]{Solution/AHSME_I_1961_40_Solution_Figure_1.png}
                    \end{figure}
                    

 From the graph, the shortest distance from the origin to the line is the \href{https://artofproblemsolving.com/wiki/index.php?title=Altitude}{altitude} to the hypotenuse of the right triangle with legs $5$ and $12$.  The hypotenuse is $13$ and the area is $30$, so the altitude to the hypotenuse is $\frac{60}{13}$, which is answer choice $\boxed{\textbf{(A)}}$.

 



\textbf{Solution 2}

Solve for $y$ in the linear equation.
 \[12y = 60 - 5x\]
 \[y = 5 - \frac{5x}{12}\]Substitute $y$ in $\sqrt{x^2+y^2}$.
 \[\sqrt{x^2 + (5 - \frac{5x}{12})^2}\]
 \[\sqrt{x^2 + 25 - \frac{25x}{6} + \frac{25x^2}{144}}\]
 \[\sqrt{\frac{169x^2}{144} - \frac{25x}{6} + 25}\]To find the minimum, find the vertex of the quadratic.  The x-value of the vertex is $\frac{25}{6} \cdot \frac{72}{169} = \frac{300}{169}$.  Thus, the minimum value is
 \[\sqrt{\frac{169}{144} \cdot \frac{300^2}{169^2} - \frac{25}{6} \cdot \frac{300}{169} + 25}\]
 \[\sqrt{\frac{10000}{16 \cdot 169} - \frac{1250}{169} + 25}\]
 \[\sqrt{\frac{625}{169} - \frac{1250}{169} + \frac{4225}{169}}\]
 \[\sqrt{\frac{3600}{169}}\]
 \[\frac{60}{13}\]The answer is $\boxed{\textbf{(A)}}$.

 



\textbf{Solution 3}

By Cauchy-Schwarz, 
 \[(5^2 + 12^2)(x^2 + y^2) \geq (5x+12y)^2\]Therefore: 
 \[(13^2)(x^2 + y^2) \geq 60^2\]
 \[13\sqrt{x^2 + y^2} \geq 60\]Therefore:
 \[\sqrt{x^2 + y^2} \geq \frac{60}{13}\]Thus the answer is $\boxed{\textbf{(A)}}$.

 


\end{enumerate}